% ============================================================
%  Memoria del Proyecto Final — Interacción Persona-Ordenador
%  Fieldnotes: Enciclopedia de Técnicas de Usabilidad
% ============================================================
\documentclass[12pt, a4paper]{report}

% ── Idioma y codificación ──────────────────────────────────
\usepackage[spanish]{babel}
\usepackage[utf8]{inputenc}
\usepackage[T1]{fontenc}

% ── Tipografía ────────────────────────────────────────────
\usepackage{lmodern}

% ── Márgenes ─────────────────────────────────────────────
\usepackage[top=2.5cm, bottom=2.5cm, left=3cm, right=2.5cm]{geometry}

% ── Interlineado ─────────────────────────────────────────
\usepackage{setspace}
\onehalfspacing

% ── Imágenes ─────────────────────────────────────────────
\usepackage{graphicx}
\usepackage{float}

% ── Colores y tablas ─────────────────────────────────────
\usepackage{xcolor}
\usepackage{booktabs}
\usepackage{longtable}
\usepackage{array}
\usepackage{tabularx}
\usepackage{multirow}

% ── Listas y enumeraciones ───────────────────────────────
\usepackage{enumitem}

% ── Hipervínculos ────────────────────────────────────────
\usepackage[hidelinks]{hyperref}

% ── Cabeceras y pies de página ───────────────────────────
\usepackage{fancyhdr}
\pagestyle{fancy}
\fancyhf{}
\fancyhead[L]{\small Proyecto Final IPO — Fieldnotes}
\fancyhead[R]{\small \leftmark}
\fancyfoot[C]{\thepage}
\renewcommand{\headrulewidth}{0.4pt}

% ── Caja de nota de ejecución ────────────────────────────
\usepackage{tcolorbox}
\tcbuselibrary{skins}
\newtcolorbox{notaejecucion}[1][]{
  colback=gray!8,
  colframe=gray!40,
  fonttitle=\bfseries\small,
  title=Contexto de ejecución,
  #1
}

% ── Placeholder de imagen ────────────────────────────────
% Uso: \imagenph{ancho}{descripción}
\newcommand{\imagenph}[2]{%
  \begin{figure}[H]
    \centering
    \fbox{\parbox[c][6cm][c]{#1}{\centering\color{gray}\textit{[Imagen: #2]}}}
    \caption{#2}
  \end{figure}%
}

% ─────────────────────────────────────────────────────────
\begin{document}

% ══════════════════════════════════════════════════════════
%  PORTADA
% ══════════════════════════════════════════════════════════
\begin{titlepage}
  \centering
  \vspace*{2cm}
  {\LARGE\bfseries Proyecto Final — Interacción Persona-Ordenador\par}
  \vspace{0.5cm}
  {\large\itshape Fieldnotes: Enciclopedia de Técnicas de Usabilidad\par}
  \vspace{3cm}
  \imagenph{8cm}{Logo o imagen de portada}
  \vfill
  {\large
    Olivia Troitiño · Víctor Gonzalo · Rodrigo Coronel \\
    Nube Arsuaga · Xianlong Wu · Eloy Sánchez-Viedma\par}
  \vspace{1cm}
  {\large Curso 2024--2025\par}
\end{titlepage}

% ══════════════════════════════════════════════════════════
%  ACLARACIONES
% ══════════════════════════════════════════════════════════
\chapter*{Aclaraciones}
\addcontentsline{toc}{chapter}{Aclaraciones}

Esta memoria documenta el proceso de diseño del proyecto final de la asignatura
\textit{Interacción Persona-Ordenador}. El objetivo del proyecto, tal como define el
enunciado, es elaborar una página web que explique técnicas de usabilidad a un perfil de
usuario concreto: \textbf{desarrolladores de software} que conocen la existencia de dichas
técnicas pero no tienen experiencia aplicándolas. El contexto de uso prevé tanto consultas
desde el puesto de oficina (durante la planificación de pruebas) como consultas rápidas
desde dispositivos móviles o tablets (durante la ejecución de las pruebas).

La memoria combina la explicación de las técnicas de HCI aplicadas y el razonamiento detrás
de cada decisión de diseño. Las fuentes utilizadas son los apuntes de clase proporcionados
por el profesor, los elaborados en clase por los alumnos, y fuentes externas. El enunciado
completo del proyecto se reproduce en el Anexo~\ref{anx:enunciado}.

% ══════════════════════════════════════════════════════════
%  TABLA DE CONTENIDOS
% ══════════════════════════════════════════════════════════
\tableofcontents
\clearpage

% ══════════════════════════════════════════════════════════
%  FASES DE DISEÑO
% ══════════════════════════════════════════════════════════
\chapter{Análisis y Planificación}

\begin{notaejecucion}
\textbf{28 de abril – 4 de mayo.} Olivia, como gestora del equipo, planifica la práctica:
estructura de la memoria, diagrama de Gantt preliminar y plan básico de fases. Comunica al
grupo que cada miembro debe ir pensando qué técnicas de usabilidad podrían implementarse
en la primera fase.

\textbf{5 de mayo.} Sesión en clase (Olivia, Víctor, Rodrigo). Se dedica tiempo a buscar webs
de referencia y anotar decisiones de diseño observadas. Olivia propone coordinar una
reunión para realizar el análisis heurístico y el focus group de Personas/Scenarios. Víctor
sugiere usar \textit{Lyssna} (Usability Hub) como herramienta de apoyo.

\textbf{8 de mayo.} Cada miembro realiza su análisis heurístico individual sobre las tres webs
de referencia, pero no se llega a discutirlo en grupo.

\textbf{12 de mayo.} Sesión en clase. Discusión grupal sobre el análisis heurístico. Se
realiza el focus group para evaluar qué decisiones de diseño implementar en nuestra web.
\end{notaejecucion}

En la fase de análisis se definen los objetivos de la página web, las características de los
usuarios objetivo y los requisitos técnicos del desarrollo.

% ──────────────────────────────────────────────────────────
\section{Contenido}

El contenido está definido como una enciclopedia dedicada a las técnicas de usabilidad de
la interacción persona-ordenador. El tiempo de vida del producto es largo, ya que se trata de
una enciclopedia o base de datos actualizable. La gestión de las actualizaciones queda en
manos del cliente, al ser un entregable con fecha límite de trabajo.

% ──────────────────────────────────────────────────────────
\section{Usuario}

La web se dirige a un usuario específico: desarrolladores de software cuyo conocimiento de
las técnicas de usabilidad está limitado al ámbito teórico, sin experiencia aplicándolas. La
aplicación se plantea como una web de difusión general —abierta al público— pero
orientada a usuarios especializados en el tema de usabilidad e interacción persona-ordenador.

% ──────────────────────────────────────────────────────────
\section{Producción}

El proyecto se organiza con los siguientes roles de equipo:
\begin{itemize}
  \item \textbf{Gestor de producción y recursos}: organización de tareas, tiempos y equipo.
  \item \textbf{Programadores y técnicos}: implementación de la web.
  \item \textbf{Documentación y guionistas}: determinan el contenido y su estructura.
\end{itemize}

El equipo está compuesto por seis participantes, todos estudiantes de un doble grado en
ingeniería informática e ingeniería en robótica. Sus perfiles son similares, aunque con
distintos niveles de experiencia en ingeniería.

% ──────────────────────────────────────────────────────────
\section{Requisitos}

Los requisitos definidos por el cliente son:
\begin{enumerate}
  \item \textbf{Usuario específico}: Desarrolladores de software que conocen la existencia de
    técnicas de usabilidad pero no tienen experiencia aplicándolas.
  \item \textbf{Objetivo específico}: Encontrar técnicas que les permitan tomar decisiones
    centradas en el usuario o evaluar sus desarrollos/prototipos, y contar con la información
    suficiente para aprender a ejecutarlas.
  \item \textbf{Contexto de uso}: Los usuarios deben poder consultar la información de
    manera eficiente desde cualquier dispositivo capaz de mostrar contenido web —tanto desde
    el puesto de oficina como desde el móvil o tablet durante la ejecución de una prueba.
\end{enumerate}

% ──────────────────────────────────────────────────────────
\section{Análisis de productos existentes}
\label{sec:analisis-productos}

Analizamos ejemplos existentes de páginas que recopilan técnicas de usabilidad y diseño.
No todas se centran exclusivamente en usabilidad ni contienen únicamente clasificación de
técnicas; aun así, se incluyeron porque ofrecen aprendizajes relevantes.

Utilizamos la \textbf{Evaluación Heurística} basada en las 10 reglas de usabilidad web de
Jakob Nielsen. Cada miembro del equipo realizó el análisis de forma independiente, y
posteriormente los resultados se pusieron en común mediante un focus group el 12 de mayo.
Este proceso combinado —análisis individual seguido de debate grupal— nos permitió
detectar tanto coincidencias claras como discrepancias que enriquecieron las conclusiones.

% ··············································
\subsection{Usability Body of Knowledge (UBK)}

\imagenph{\linewidth}{Captura de pantalla de UBK}

Los análisis individuales del equipo coincidieron en los siguientes puntos sobre UBK:

\begin{itemize}
  \item \textbf{H1 – Visibilidad del estado} (\checkmark): La ruta de navegación
    (\textit{breadcrumbs}) indica claramente al usuario en qué sección se ubica
    (ej.\ \texttt{Home > Topics > Methods > User Research}).
  \item \textbf{H2 – Lenguaje del usuario} (\textasciitilde): El lenguaje es familiar, pero la
    organización de la página resulta por momentos inconsistente.
  \item \textbf{H3 – Control y libertad} (\checkmark): El usuario puede volver atrás en la
    jerarquía haciendo clic en los enlaces de la ruta superior.
  \item \textbf{H4 – Consistencia y estándares} (\textasciitilde): Usa signos \texttt{+}/\texttt{-}
    para desplegar subsecciones del menú, estándar de hace veinte años.
  \item \textbf{H5 – Prevención de errores} (—): Al ser un sitio casi estático, el riesgo de
    error de usuario es mínimo, pero no hay elementos activos de prevención.
  \item \textbf{H6 – Reconocimiento sobre memorización} (\textasciitilde): Las categorías
    visibles ayudan, pero no hay elementos visuales de apoyo.
  \item \textbf{H7 – Flexibilidad y eficiencia} (\textasciitilde): Eficiente para el experto que
    busca un dato; más lento para el aprendizaje.
  \item \textbf{H8 – Diseño estético y minimalista} (\texttimes): El diseño es anticuado y la
    falta de espacio en blanco (\textit{whitespace}) fatiga la vista.
  \item \textbf{H9 – Reconocimiento y recuperación de errores} (\checkmark): Si el usuario
    busca una técnica inexistente, la web le avisa que no hay resultados.
  \item \textbf{H10 – Ayuda y documentación} (\texttimes): No hay FAQs ni documentación
    de uso de la propia web.
\end{itemize}

\textbf{Notas del focus group (12 de mayo):}
\begin{itemize}
  \item Demasiadas sub-secciones; la navegación no es intuitiva.
  \item Podrían usarse más colores para indicar estado y secciones.
  \item Hay enlaces desactualizados que no llevan a ningún lado.
  \item Las descripciones de las técnicas podrían estar mejor organizadas y son inconsistentes entre sí.
\end{itemize}

% ··············································
\subsection{No Solo Usabilidad (NSU)}

\imagenph{\linewidth}{Captura de pantalla de NSU}

NSU funciona más como una revista/blog que como una enciclopedia de técnicas, pero
resultó útil para observar cómo se describe cada técnica en un formato editorial. El equipo
destacó:

\begin{itemize}
  \item \textbf{H1–H4} (\checkmark / \textasciitilde): La navegación es clara y el lenguaje
    profesional y cercano; algunos artículos presentan estructuras visuales diferentes.
  \item \textbf{H5 – Prevención de errores} (\texttimes): La barra de búsqueda dice
    únicamente \textit{"powered by Google"} sin explicar qué se está buscando.
  \item \textbf{H6 – Reconocimiento sobre memorización} (\texttimes): Al hacer scroll el
    menú de navegación desaparece, obligando al usuario a memorizar las secciones.
  \item \textbf{H8 – Diseño estético} (\textasciitilde): Las imágenes cubren un tercio de la
    página, entorpeciendo la navegación y distrayendo de la lectura.
  \item \textbf{H10 – Ayuda y documentación} (\texttimes): Sin FAQs ni manual de uso.
\end{itemize}

\textbf{Notas del focus group:}
\begin{itemize}
  \item Lenguaje profesional y legible.
  \item Faltan iconos que faciliten la navegación y la comprensión de la arquitectura de información.
  \item La web funciona como página lineal de scroll-down: buena organización del contenido,
    pero estructura de acceso poco clara.
\end{itemize}

% ··············································
\subsection{Service Design Tools (SDT)}

\imagenph{\linewidth}{Captura de pantalla de SDT}

SDT destaca como la referencia más moderna de las tres. Permite clasificar las técnicas
por \textit{When} (etapa del desarrollo), \textit{Who} (a quién involucrar), \textit{What} (qué
aspecto trabajar) y \textit{How} (tipo de representación).

\begin{itemize}
  \item \textbf{H1 – Visibilidad del estado} (\checkmark): Feedback visual inmediato al
    seleccionar filtros; la sección activa se subraya en el menú.
  \item \textbf{H3 – Control y libertad} (\checkmark): Una \texttt{X} permite quitar cualquier
    filtro sin recargar.
  \item \textbf{H5 – Prevención de errores} (\checkmark): El buscador tiene un
    \textit{placeholder} que aclara qué se puede buscar; el sistema de filtrado evita
    llegar a páginas vacías.
  \item \textbf{H6 – Reconocimiento sobre memorización} (\checkmark): Los filtros aplicados
    permanecen visibles en una lista persistente.
  \item \textbf{H8 – Diseño estético} (\checkmark): La mejor de las tres en este apartado:
    minimalista, con buen uso del espacio en blanco y tipografía legible.
  \item \textbf{H10 – Ayuda y documentación} (\textasciitilde): Tiene sección de tutoriales,
    pero falta una guía de inicio rápido.
\end{itemize}

\textbf{Notas del focus group:}
\begin{itemize}
  \item Organización lógica y simple; buen uso de iconos; estética bien lograda e intuitiva.
  \item Los filtros pueden llegar a confundir si se combinan demasiados.
  \item La descripción de las técnicas es más informativa que orientada a la ejecución;
    solo en algunos casos incluye el \textit{cómo llevarlo a cabo}.
\end{itemize}

% ··············································
\subsection{Conclusiones del análisis}
\label{sec:conclusiones-analisis}

\begin{itemize}
  \item SDT es la mejor de las tres en cuanto a usabilidad y diseño visual.
  \item UBK y SDT describen las técnicas con más rigor, pero resultan demasiado extensas y
    difíciles de leer para nuestro caso de uso: consulta durante la ejecución de una prueba.
  \item Ninguna de las tres ofrece FAQs ni documentación de ayuda sobre la propia web.
  \item La mejor opción sería una mezcla: la estructura de acceso y la estética de SDT, la
    profundidad de contenido de UBK, y la accesibilidad de lenguaje de NSU.
\end{itemize}

Estas conclusiones, junto con las notas del Affinity Diagram (véase el capítulo~\ref{ch:diseno}),
guiaron directamente las decisiones de diseño del prototipo.

% ══════════════════════════════════════════════════════════
\chapter{Modelado del Usuario}

\begin{notaejecucion}
\textbf{12 de mayo.} Al concluir el focus group del análisis heurístico, se decide comenzar
el modelado del usuario. Cada miembro del equipo debe crear una o más Persona/Scenario
para el 13 de mayo.

\textbf{22 de mayo.} Cada miembro revisa su Persona/Scenario y se crean los casos de uso
a partir de los perfiles definidos.
\end{notaejecucion}

La fase de modelado del usuario se caracteriza por la definición de los perfiles de los
usuarios de la aplicación. Se divide en dos partes: \textbf{Personas y Scenarios}, y
\textbf{Casos de Uso}.

% ──────────────────────────────────────────────────────────
\section{Personas y Scenarios}

Las Personas son usuarios realistas con características determinadas, descritos de manera
narrativa con identidades concretas. Los Scenarios son casos específicos de utilización,
teniendo en cuenta las tareas que el sistema debe llevar a cabo y el contexto en que la
persona interactuará con la aplicación.\footnote{Toni Granollers, Jesús Lorés y José J.\ Cañas
— \textit{Diseño de sistemas interactivos centrados en el usuario} (UOC, 2005).}

Los tipos de persona definidos para este proyecto son:

\begin{description}
  \item[Focal] Principal tipo de usuario al que se dirige el proyecto: desarrolladores con
    conocimiento teórico de las técnicas de usabilidad pero sin experiencia aplicándolas.
  \item[Secundaria] Usan el producto de manera satisfactoria; en caso de conflicto de
    intereses, se prioriza al usuario Focal. En este proyecto: desarrolladores sin
    conocimiento previo, personas con interés general en usabilidad, y gestores o jefes
    que quieren dirigir a sus equipos de diseño.
  \item[No prioritario] Usan el producto de manera infrecuente; baja prioridad.
  \item[Involucrado] No usan el producto pero se ven afectados por sus resultados (ej.\
    participantes en pruebas de usabilidad conducidas por los desarrolladores que aprenden
    en nuestra web).
  \item[Excluido] No usan el producto ni se ven afectados.
\end{description}

La descripción de cada persona incluye un perfil geográfico, demográfico y psicosocial, así
como su relación con el producto (rol, frecuencia de uso, porcentaje del uso total). Los
scenarios recogen el contexto de utilización, el desarrollo de la interacción y las
características emocionales.

% ············· FOCAL ·································
\subsection{Usuarios Focales}

% --- Marcos ---
\subsubsection{Marcos}

\begin{table}[H]
\centering
\begin{tabularx}{\linewidth}{lX}
\toprule
\textbf{Edad} & 29 años \\
\textbf{Profesión} & Desarrollador full-stack \\
\midrule
\textbf{Descripción} &
Marcos vive en Buenos Aires. Estudió Ingeniería en Informática en la UBA, donde aprendió
sobre técnicas de usabilidad. Lleva dos años trabajando en una startup de 15 personas como
desarrollador y fue elegido recientemente para liderar el próximo proyecto. Al ser una empresa
pequeña, no hay diseñador UX ni fondos para investigación de usuario, por lo que esta
responsabilidad recae en Marcos. Tiene una laptop para desarrollo y una tablet para organización.
Le gusta el contenido directo con ejemplos, y es impaciente con textos extensos y webs difíciles
de navegar. \\[4pt]
\textbf{Scenario} &
Es lunes por la mañana. El equipo ha terminado un prototipo de la pantalla principal y Marcos
necesita validar con usuarios que el diseño se entiende —tiene hasta el viernes. Abre la laptop,
busca "técnicas de usabilidad", entra en la web y aplica filtros (bajo presupuesto, prototipo, poco
tiempo). Encuentra dos opciones, abre sus fichas y se queda con la que tiene los pasos más claros
y un ejemplo concreto. El jueves, mientras ejecuta la prueba, usa la tablet para consultar
rápidamente un paso de la técnica que no recordaba. \\
\bottomrule
\end{tabularx}
\caption{Persona Focal — Marcos}
\end{table}

% --- Alejandro ---
\subsubsection{Alejandro}

\begin{table}[H]
\centering
\begin{tabularx}{\linewidth}{lX}
\toprule
\textbf{Edad} & 24 años \\
\textbf{Profesión} & Estudiante de ingeniería informática y desarrollador web junior \\
\midrule
\textbf{Descripción} &
Alejandro vive en Madrid y combina estudios universitarios con pequeños proyectos de desarrollo
web. Tiene conocimientos generales sobre experiencia de usuario y sabe que existen técnicas de
usabilidad. Accede a internet principalmente desde su portátil. Pierde interés rápidamente cuando
una web tiene demasiado texto, navegación confusa o información mal organizada. \\[4pt]
\textbf{Scenario} &
Es un miércoles por la tarde. Debe preparar el diseño inicial de una aplicación para una práctica
universitaria y el profesor le pide justificar las técnicas de usabilidad utilizadas. Alejandro prueba
primero una web con demasiados bloques de texto y se frustra. Luego accede a otra web con
clasificación por categorías y filtros visuales y localiza rápidamente la información que necesita. \\
\bottomrule
\end{tabularx}
\caption{Persona Focal — Alejandro}
\end{table}

% --- Elena ---
\subsubsection{Elena}

\begin{table}[H]
\centering
\begin{tabularx}{\linewidth}{lX}
\toprule
\textbf{Edad} & 34 años \\
\textbf{Profesión} & Diseñadora de productos en una agencia de marketing digital \\
\midrule
\textbf{Descripción} &
Elena vive en Madrid y lleva 8 años en el mundo del desarrollo web. Su rol es híbrido: aunque su
fuerte es programar, ahora también asume responsabilidades de UX con clientes exigentes. Conoce
la teoría de usabilidad, pero trabaja a un ritmo frenético y necesita adaptar metodologías complejas
a formato ágil. Usuaria Focal con frecuencia de uso muy alta e interacción intensiva al arrancar cada
proyecto. \\[4pt]
\textbf{Scenario} &
Es miércoles por la tarde. Un cliente le presiona porque desconfía del nuevo flujo de pago de su
e-commerce; el objetivo SMART de Elena es obtener datos de usuarios en menos de 48 horas
sin coste. Filtra la web por "menos de 2 horas de preparación", encuentra un test remoto, lee la
ficha al detalle y descarga una plantilla de guion que personaliza de inmediato. \\
\bottomrule
\end{tabularx}
\caption{Persona Focal — Elena}
\end{table}

% --- Li Xiaoxiao ---
\subsubsection{Li Xiaoxiao}

\begin{table}[H]
\centering
\begin{tabularx}{\linewidth}{lX}
\toprule
\textbf{Edad} & 21 años \\
\textbf{Profesión} & Estudiante de segundo curso, Universidad Tsinghua \\
\midrule
\textbf{Descripción} &
Li Xiaoxiao cursa Ingeniería de Software y Diseño de Interacción. Su meta es realizar un máster en
UX/Usabilidad en España. Tiene conocimientos teóricos pero sin experiencia práctica. Necesita
guías metodológicas sencillas, paso a paso, en un lenguaje accesible (el vocabulario técnico en
español a veces le resulta complejo). \\[4pt]
\textbf{Scenario} &
Es jueves a las 19:00 en una cafetería de Barcelona. Xiaoxiao necesita seleccionar una técnica de
usabilidad para evaluar su prototipo (entrega mañana). Accede desde su ordenador a la web,
busca una ficha técnica clara que explique materiales, pasos y ejemplos para ejecutar la técnica
esa misma noche. La cafetería cierra pronto y está cansada tras un día de turismo y visitas
académicas. \\
\bottomrule
\end{tabularx}
\caption{Persona Focal — Li Xiaoxiao}
\end{table}

% ············· SECUNDARIA ·····························
\subsection{Usuarios Secundarios}

% --- Giovanna ---
\subsubsection{Giovanna}

\begin{table}[H]
\centering
\begin{tabularx}{\linewidth}{lX}
\toprule
\textbf{Edad} & 27 años \\
\textbf{Profesión} & Propietaria de un club de entretenimiento nocturno \\
\midrule
\textbf{Descripción} &
Giovanna no tiene tiempo para manuales densos; busca herramientas que le permitan optimizar
la tasa de conversión de reservas sin sacrificar estética. Maneja su negocio desde móvil y laptop.
Usuaria Secundaria con uso ocasional. \\[4pt]
\textbf{Scenario} &
Es medianoche. Observa que el proceso de reserva de su nuevo sistema de tablets resulta
"burocrático". Busca en la web una técnica centrada en la estética, filtra por "análisis visual" y
guarda la ficha para reunirse con sus desarrolladores al día siguiente. \\
\bottomrule
\end{tabularx}
\caption{Persona Secundaria — Giovanna}
\end{table}

% --- Lucía ---
\subsubsection{Lucía}

\begin{table}[H]
\centering
\begin{tabularx}{\linewidth}{lX}
\toprule
\textbf{Edad} & 23 años \\
\textbf{Profesión} & Estudiante de último año de Ingeniería Informática \\
\midrule
\textbf{Descripción} &
Lucía nunca ha cursado una asignatura de usabilidad; términos como "test de usabilidad" o
"evaluación heurística" le suenan vagamente. Tiene paciencia para leer si el contenido le aporta,
pero se pierde cuando un texto da por hecho conocimientos que ella no tiene. Trabaja desde portátil
y hace consultas rápidas desde el móvil. \\[4pt]
\textbf{Scenario} &
Es un miércoles. Su tutor del TFG le pide que justifique decisiones de diseño y haga alguna prueba
con usuarios. Busca "cómo evaluar la usabilidad de una web", llega a la web, entra en la sección
de introducción para orientarse, sigue un enlace a una técnica sencilla y guarda el enlace para
consultarlo desde el móvil en el metro. \\
\bottomrule
\end{tabularx}
\caption{Persona Secundaria — Lucía}
\end{table}

% --- Javier ---
\subsubsection{Javier}

\begin{table}[H]
\centering
\begin{tabularx}{\linewidth}{lX}
\toprule
\textbf{Edad} & 45 años \\
\textbf{Profesión} & Director de marketing en empresa mediana de retail \\
\midrule
\textbf{Descripción} &
Perfil puramente de negocio: no sabe programar ni diseñar, pero decide los presupuestos de
proyectos digitales. Ve la usabilidad como herramienta para vender más. A veces se siente
intimidado por la jerga técnica. Usuario Secundario con uso bajo y ocasional. \\[4pt]
\textbf{Scenario} &
Las métricas muestran que el tráfico es alto pero las ventas han caído tras la última actualización.
Javier busca una técnica que pueda sugerir a su equipo técnico sin cuestionar su trabajo. Encuentra
un artículo sobre testeo con usuarios enfocado a negocio, con ejemplos claros de éxito. \\
\bottomrule
\end{tabularx}
\caption{Persona Secundaria — Javier}
\end{table}

% --- Laura ---
\subsubsection{Laura}

\begin{table}[H]
\centering
\begin{tabularx}{\linewidth}{lX}
\toprule
\textbf{Edad} & 22 años \\
\textbf{Profesión} & Perfil de ciberseguridad ofensiva (ficticio) \\
\midrule
\textbf{Descripción} &
Persona de perfil extremo incluida para explorar cómo la web sería usada —o usada mal— bajo
un contexto de alta presión y motivos no convencionales. Accede desde dispositivos en entornos
de alta distracción. \\[4pt]
\textbf{Scenario} &
Necesita entender cómo un usuario común interactúa con una interfaz bajo estrés. Lee la ficha
del Recorrido Cognitivo (\textit{Cognitive Walkthrough}) desde el móvil para identificar en qué pasos
el flujo genera más confusión. \\
\bottomrule
\end{tabularx}
\caption{Persona Secundaria — Laura}
\end{table}

% ············· INVOLUCRADO ····························
\subsection{Usuario Involucrado}

\subsubsection{Sofía}

\begin{table}[H]
\centering
\begin{tabularx}{\linewidth}{lX}
\toprule
\textbf{Edad} & 26 años \\
\textbf{Profesión} & Administrativa en Recursos Humanos \\
\midrule
\textbf{Descripción} &
Sofía nunca entrará a la web de técnicas ni consumirá su contenido, pero el éxito de las técnicas
que los desarrolladores aprendan ahí afectará directamente a su experiencia como participante en
pruebas. Es usuaria de las aplicaciones que los desarrolladores evalúan. Se frustra si una interfaz
la hace sentir poco capaz. \\[4pt]
\textbf{Scenario} &
Un desarrollador le pide que use un prototipo durante diez minutos mientras él observa. Sofía
arranca el test un poco nerviosa y cohibida; espera una experiencia clara que no la haga perder
el tiempo ni sentirse juzgada. \\
\bottomrule
\end{tabularx}
\caption{Persona Involucrada — Sofía}
\end{table}

% ············· NO PRIORITARIO ·························
\subsection{Usuario No Prioritario}

\subsubsection{Carlos}

\begin{table}[H]
\centering
\begin{tabularx}{\linewidth}{lX}
\toprule
\textbf{Edad} & 40 años \\
\textbf{Profesión} & Profesor de bachillerato \\
\midrule
\textbf{Descripción} & Carlos da clase de tecnología. Oye a un alumno mencionar "pruebas de usabilidad" y siente curiosidad. \\[4pt]
\textbf{Scenario} & Busca el término en Google, entra en la web, explora la página de inicio y las FAQs durante unos diez minutos y luego sigue con su vida. \\
\bottomrule
\end{tabularx}
\caption{Persona No Prioritaria — Carlos}
\end{table}

% ──────────────────────────────────────────────────────────
\section{Casos de Uso}

A partir de los resultados de Personas/Scenarios se crean los casos de uso. Un caso de uso
describe cómo un actor interactúa con el sistema para alcanzar un objetivo concreto: qué
necesita, qué pasos sigue, y qué puede salir bien o mal. No todos los casos de uso se
implementarán necesariamente, pero sus descripciones orientan las decisiones de diseño.

% --- CU1 ---
\subsection{CU1: Consultar la ficha de una técnica concreta para ejecutarla}

\begin{longtable}{lp{10cm}}
\toprule
\textbf{Objetivo} & El desarrollador accede a toda la información necesaria de una técnica para poder aplicarla en su propio proyecto. \\
\textbf{Precondiciones} & Usuario con acceso a internet y a un navegador. La web está disponible. El usuario tiene alguna necesidad relacionada con la usabilidad, aunque no necesariamente sabe qué técnica concreta la resolverá. \\
\textbf{Post-condiciones} & El sistema ha mostrado al usuario el contenido de una o varias técnicas. \\
\textbf{Cond. final éxito} & El usuario ha localizado la técnica adecuada y dispone de información suficiente para ejecutarla. \\
\textbf{Cond. final fallo} & El usuario no encuentra ninguna técnica adecuada, o la ficha no contiene información suficiente. \\
\textbf{Actor primario} & Desarrollador de software sin experiencia previa en técnicas de usabilidad. \\
\midrule
\textbf{Secuencia normal} & 1. Entra en la web. \newline 2. Explora el listado/categorías. \newline 3. Aplica filtros (fase, objetivo, tipo). \newline 4. Selecciona una técnica. \newline 5. Lee la ficha. \\
\textbf{Alternativa a} & 2a. Usa el buscador por nombre. \newline 3a. Selecciona la técnica de los resultados. \newline 4a. Lee la ficha. \\
\textbf{Alternativa b} & 2b. Accede a una sección de ayuda/recomendaciones. \newline 3b. Sigue una sugerencia que enlaza directamente a una técnica. \newline 4b. Lee la ficha. \\
\bottomrule
\end{longtable}

% --- CU2 ---
\subsection{CU2: Guardar técnicas de usabilidad para consultas posteriores}

\begin{longtable}{lp{10cm}}
\toprule
\textbf{Objetivo} & Almacenar técnicas relevantes para recuperarlas más adelante durante el desarrollo de un proyecto o la preparación de documentación. \\
\textbf{Actor primario} & Estudiante de ingeniería informática y desarrollador web junior. \\
\textbf{Actores secundarios} & Profesor o tutor académico que solicita justificar las técnicas utilizadas. \\
\midrule
\textbf{Secuencia normal} &
1. Navega por las técnicas disponibles. \newline
2. Consulta distintas fichas relacionadas con su proyecto. \newline
3. Marca una técnica que considera útil. \newline
4. El sistema almacena la selección. \newline
5. Continúa explorando. \newline
6. Añade otras técnicas de interés. \newline
7. Revisa el listado de técnicas guardadas. \\
\bottomrule
\end{longtable}

% --- CU3 ---
\subsection{CU3: Diagnóstico y prescripción de técnicas de evaluación}

\begin{longtable}{lp{10cm}}
\toprule
\textbf{Objetivo} & Guiar al diseñador en la identificación, selección e instrucción de la metodología de evaluación óptima en función de las restricciones de su proyecto. \\
\textbf{Actor primario} & Diseñador/evaluador novel que requiere validación metodológica externa. \\
\midrule
\textbf{Secuencia normal} &
1. Accede al módulo de diagnóstico. \newline
2. Introduce los síntomas observados en su interfaz mediante un menú de opciones múltiples. \newline
3. Especifica los recursos críticos disponibles (tiempo, participantes, presupuesto). \newline
4. El sistema calcula la compatibilidad y presenta la técnica ideal recomendada. \newline
5. El usuario confirma la selección y el sistema despliega el plan de ejecución estructurado. \\
\bottomrule
\end{longtable}

% --- CU4 ---
\subsection{CU4: Calcular el número de usuarios necesarios y generar mensaje de reclutamiento}

\begin{longtable}{lp{10cm}}
\toprule
\textbf{Objetivo} & Que el desarrollador sepa cuántos usuarios necesita para que su prueba sea válida y obtenga un texto listo para reclutar participantes. \\
\textbf{Actor primario} & Desarrollador de software sin experiencia previa. \\
\midrule
\textbf{Secuencia normal} &
1. Entra en "Calculadora de Usuarios y Reclutamiento". \newline
2. Selecciona el tipo de público de su app (general / varios perfiles / muy específico). \newline
3. Introduce el tipo de técnica mediante un menú desplegable. \newline
4. El sistema aplica la fórmula y muestra el número recomendado de participantes. \newline
5. El sistema genera una plantilla de mensaje de reclutamiento personalizada. \newline
6. El usuario pulsa "Copiar texto". \\
\bottomrule
\end{longtable}

% --- CU5 ---
\subsection{CU5: Descargar recursos y plantillas para la ejecución ágil de una técnica}

\begin{longtable}{lp{10cm}}
\toprule
\textbf{Objetivo} & Acceder a herramientas listas para usar (plantillas de guiones, checklists, matrices) asociadas a una técnica para aplicarla de forma inmediata. \\
\textbf{Actor primario} & Usuario Focal (desarrollador con conocimientos en usabilidad). \\
\textbf{Actores secundarios} & Usuario Secundario con conocimientos previos que requiere herramientas ejecutables. \\
\midrule
\textbf{Secuencia normal} &
1. Entra en la web. \newline
2. Busca la técnica mediante filtros o el menú. \newline
3. Selecciona la técnica. \newline
4. Accede a la sección "Recursos y Plantillas" dentro de la ficha. \newline
5. Selecciona el archivo que necesita. \newline
6. El sistema inicia la descarga. \\
\bottomrule
\end{longtable}

% --- CU6 ---
\subsection{CU6: Explorar la web de forma casual}

\begin{longtable}{lp{10cm}}
\toprule
\textbf{Objetivo} & Que un usuario sin perfil técnico pueda entender de forma básica qué es una técnica de usabilidad navegando por curiosidad, sin objetivo profesional concreto. \\
\textbf{Actor primario} & Usuario no prioritario (sin perfil técnico, curiosidad puntual). \\
\midrule
\textbf{Secuencia normal} &
1. Busca "test de usabilidad" en su navegador. \newline
2. Entra en la web. \newline
3. Explora la página de inicio y las FAQs. \newline
4. Lee acerca de alguna técnica concreta. \\
\textbf{Cond. final fallo} & El usuario encuentra el contenido demasiado técnico y abandona la web. \\
\bottomrule
\end{longtable}

% ══════════════════════════════════════════════════════════
\chapter{Diseño}
\label{ch:diseno}

\begin{notaejecucion}
\textbf{24 de mayo.} Se realiza el Affinity Diagram de forma colaborativa y remota. A
continuación se celebra un focus group para discutir qué grupos y filtros debían mantenerse,
y cuáles fusionarse o eliminarse.
\end{notaejecucion}

La fase de diseño responde a las características definidas en el análisis y el modelado del
usuario. Su objetivo es traducir los hallazgos de las fases anteriores en decisiones concretas
sobre la estructura y el contenido de la web. La técnica principal empleada fue el
\textbf{Affinity Diagram}.

% ──────────────────────────────────────────────────────────
\section{Affinity Diagram}

El Affinity Diagram es una técnica de síntesis que permite organizar grandes cantidades de
datos cualitativos en grupos temáticos emergentes. Se aplica en tres pasos: (1) generación
individual de notas/observaciones, (2) agrupación silenciosa y colaborativa, y (3) debate y
nombramiento de los grupos resultantes. Elegimos esta técnica porque en las fases
anteriores habíamos acumulado una gran cantidad de notas heterogéneas —del análisis
heurístico, los focus groups y los casos de uso— que necesitaban sintetizarse antes de
diseñar los wireframes.

\subsection{Paso 1: Aportación individual de notas}

Cada miembro del equipo aportó observaciones independientes basadas en el análisis
heurístico, los scenarios y los casos de uso. A continuación se reproducen las contribuciones
individuales agrupadas por autor.

\imagenph{\linewidth}{Tablero completo con las notas individuales del Affinity Diagram (captura de pantalla de la herramienta online)}

\paragraph{Olivia Troitiño}
\begin{itemize}[nosep]
  \item Faltan íconos en las webs analizadas.
  \item Las webs en general podrían ser más concisas; pierde interés cuando hay demasiado texto.
  \item Prefiere lenguaje accesible; busca información clara porque trabaja con tiempos limitados.
  \item UBK tiene demasiadas sub-secciones y una navegación poco intuitiva; hay enlaces desactualizados.
  \item Las webs en general no tienen FAQs.
  \item La ficha no siempre contiene información suficiente para aplicar la técnica.
  \item Organización por clasificación, presupuesto, fases, resultado u objetivo.
  \item Filtros por tiempos y por tipo de análisis.
  \item El sistema debe permitir almacenar selecciones del usuario.
  \item Posibilidad de comparar dos técnicas.
\end{itemize}

\paragraph{Rodrigo Coronel}
\begin{itemize}[nosep]
  \item Las interfaces simples con filtros y clasificación visual permiten localizar contenido con
    menos esfuerzo y mejoran la accesibilidad para usuarios noveles.
  \item Los usuarios esperan experiencias directas, intuitivas y con baja carga cognitiva, especialmente
    bajo presión o tiempo limitado.
  \item Las webs analizadas presentan exceso de texto, demasiadas subsecciones y estructuras complejas.
  \item La web debe facilitar la aplicación práctica mediante fichas claras, recomendaciones, recursos
    descargables y acceso rápido.
  \item Los usuarios necesitan recuperar información previamente consultada sin repetir búsquedas.
\end{itemize}

\paragraph{Víctor Gonzalo}
\begin{itemize}[nosep]
  \item SDT tiene una iconografía bonita y convencional para sus técnicas.
  \item "Necesito que la información sea muy digerible, sin código ni tecnicismos absurdos."
  \item NSU no tiene una forma clara de acceder a las técnicas; parece una página lineal de scroll-down.
  \item UBK tiene una gran organización de la información y distintas \textit{landings} bien organizadas.
  \item SDT tiene buena estructura (when/who/what/how) pero demasiado amplia.
  \item No puede ocurrir que el enlace de descarga de recursos esté roto.
\end{itemize}

\paragraph{Eloy Sánchez-Viedma}
\begin{itemize}[nosep]
  \item La página de inicio debe invitar al usuario a explorarla, sin sobrecarga y con FAQs.
  \item La información debe explicarse de forma que cualquiera pueda entenderla.
  \item Se debe aclarar el uso de cada técnica para que el usuario encuentre la adecuada en el menor
    tiempo posible.
\end{itemize}

\paragraph{Nube Arsuaga}
\begin{itemize}[nosep]
  \item "Odio las webs que son solo texto infumable. Si me vas a explicar la usabilidad, demuéstralo
    con la usabilidad de tu propia página."
  \item Necesitan ejemplos reales de empresas que aplicaron la técnica.
  \item La web debe cargar rápido en dispositivos móviles.
  \item "Cuando tengo una duda en medio de un taller, entro en Google rápido. Necesito un buscador
    que vaya al grano, no un artículo de opinión."
  \item El usuario busca filtrar técnicas por la fase del proyecto en que se encuentra.
\end{itemize}

\paragraph{Xianlong Wu (昱龙吴)}
\begin{itemize}[nosep]
  \item UBK cumple con la visibilidad del estado pero es muy pobre visualmente.
  \item Necesidad de una ficha técnica clara con lenguaje accesible.
  \item La ausencia de menú persistente en NSU obliga a recordar dónde estaban las secciones.
  \item Evitar imágenes desproporcionadas que rompen la continuidad de la lectura.
  \item Procurar que los recursos no estén dañados (enlaces caducados, etc.).
\end{itemize}

\subsection{Paso 2: Agrupación silenciosa y colaborativa}

La agrupación se realizó de forma online y simultánea. La moderadora (Olivia) observó la
evolución del tablero sin intervenir en las decisiones de agrupación, apuntando únicamente
sus observaciones:

\begin{itemize}
  \item En los primeros diez minutos aparecieron los siguientes grupos emergentes: filtros y
    organización (que de a poco se dividió en dos), información comprensible y concisa,
    problemas a evitar + reducción de carga cognitiva (que también comenzó a separarse),
    plantillas/recursos, e iconografía y estilo visual.
  \item \textbf{Momento de desacuerdo 1}: La tarjeta "página de inicio debe ser concisa" fue
    movida repetidamente entre "Información comprensible y concisa" y "Problemas a evitar",
    lo que acabó resolviéndose en el debate de la fase 3.
  \item \textbf{Momento de desacuerdo 2}: Dos integrantes reclamaron simultáneamente la
    tarjeta "Faltan íconos"; uno la tomó y el otro aceptó el cambio.
  \item El grupo de plantillas absorbió tarjetas de recursos/descargables e imágenes, incluyendo
    "La web debe cargar rápido en dispositivos móviles", cuya conexión con el grupo no resultó
    obvia a la moderadora.
  \item También se detectó solapamiento entre los dos sub-grupos de filtros y organización, y
    entre el grupo de "Problemas a evitar" y el de "Información comprensible y concisa".
\end{itemize}

\imagenph{\linewidth}{Tablero tras la agrupación silenciosa (captura de pantalla)}

\subsection{Paso 3: Debate y nombrado de los grupos}

Tras el debate, el equipo acordó los siguientes cinco grupos temáticos:

\subsubsection{Diseño Visual}
Incorporar iconografía; diseño minimalista y simple; la página de inicio debe invitar al usuario
a explorarla.

\subsubsection{Información Clara y Concisa}
Información fácilmente digerible, sin sobrecarga; FAQs en la página de inicio; lenguaje accesible
para el usuario no técnico; ejemplos reales; fichas técnicas claras; buscador específico.

\subsubsection{Problemas a Evitar}
Exceso de texto, subsecciones y estructuras complejas; webs lentas; menú no persistente; diseños
anticuados; enlaces caducados; fichas sin suficiente información para ejecutar la técnica.

\subsubsection{Recursos / Accesibilidad}
Repositorio de plantillas accesible desde la página de inicio; fichas con plantillas descargables
y en formatos compatibles; la web debe cargar rápido en móvil; no tener imágenes desproporcionadas.

\subsubsection{Features (Organización, Filtrado y Funcionalidades)}
Organización de técnicas por clasificación, presupuesto, fases y resultado/objetivo; filtros por
tiempos y tipo de análisis; página de "favoritos"; comparación de dos técnicas; buscador general.

\imagenph{\linewidth}{Tablero final con los grupos nombrados (captura de pantalla)}

% ──────────────────────────────────────────────────────────
\section{De las conclusiones del diseño a los wireframes}

Los cinco grupos del Affinity Diagram condensaron las necesidades identificadas en el
análisis y el modelado del usuario en un conjunto de requisitos de diseño concretos. El
siguiente paso fue traducirlos en pantallas tangibles mediante wireframes. Dado que el
contexto de uso incluye explícitamente el acceso desde móvil durante la ejecución de una
prueba (requisito 3), adoptamos un enfoque \textbf{mobile-first}: diseñamos primero las
pantallas para el ancho de pantalla más pequeño y luego ampliamos el diseño a escritorio.

Las decisiones principales que guiaron los wireframes fueron:
\begin{itemize}
  \item Navegación persistente y visible en todo momento (soluciona H6 detectado en NSU y UBK).
  \item Filtros prominentes desde la pantalla de listado de técnicas (grupo "Features").
  \item Página de inicio concisa con una breve descripción y acceso directo al catálogo
    (grupos "Información clara y concisa" y "Diseño visual").
  \item Ficha de técnica con secciones claras y descarga de plantilla integrada
    (grupos "Recursos/Accesibilidad" e "Información clara y concisa").
\end{itemize}

% ══════════════════════════════════════════════════════════
\chapter{Prototipado}

\begin{notaejecucion}
\textbf{24 de mayo.} Tras el Affinity Diagram y el focus group, se diseñaron los wireframes.
Se adoptó un enfoque mobile-first y luego se extendió el diseño a escritorio.
\end{notaejecucion}

La fase de prototipado materializa las decisiones de diseño en representaciones visuales
concretas. En esta fase se elaboraron wireframes de baja-media fidelidad para las pantallas
principales de la web.

% ──────────────────────────────────────────────────────────
\section{Wireframes — Mobile-first}

\imagenph{\linewidth}{Wireframe mobile: página de inicio}
\imagenph{\linewidth}{Wireframe mobile: listado y filtrado de técnicas}
\imagenph{\linewidth}{Wireframe mobile: ficha de una técnica}

% ──────────────────────────────────────────────────────────
\section{Wireframes — Desktop}

\imagenph{\linewidth}{Wireframe desktop: página de inicio}
\imagenph{\linewidth}{Wireframe desktop: listado y filtrado de técnicas}
\imagenph{\linewidth}{Wireframe desktop: ficha de una técnica}

% ══════════════════════════════════════════════════════════
%  ANEXOS
% ══════════════════════════════════════════════════════════
\appendix

\chapter{Ficha de Técnica}
\label{anx:ficha}

La siguiente ficha es la plantilla utilizada como unidad de contenido de la web. El orden
de los componentes será ajustado durante la implementación: dado que los usuarios focales
ya tienen conocimiento previo de las técnicas, consideramos que el apartado "¿Cómo se
lleva a cabo?" resulta más relevante que los componentes puramente descriptivos.

\begin{table}[H]
\centering
\begin{tabularx}{\linewidth}{lX}
\toprule
Nombre de la técnica & \\
Clasificación & (Inspecting / Testing / Inquiring) \\
¿Para qué sirve? & \\
¿Para qué no sirve? & \\
Momento del ciclo de vida & \\
Tiempos & \\
Cantidad de usuarios o evaluadores & \\
Materiales & \\
Costo monetario & \\
¿Cómo se lleva a cabo? & \\
Resultado/Output & \\
Dimensiones de usabilidad evaluadas & \\
\bottomrule
\end{tabularx}
\caption{Plantilla de la ficha de técnica}
\end{table}

% ──────────────────────────────────────────────────────────
\chapter{Templates}
\label{anx:templates}

\section{Plantilla Persona-Scenario}

\begin{table}[H]
\centering
\begin{tabularx}{\linewidth}{lX}
\toprule
Edad & \\
Profesión & \\
Descripción de la persona & \\
Descripción del scenario & \\
\bottomrule
\end{tabularx}
\caption{Plantilla de Persona-Scenario}
\end{table}

\section{Plantilla Caso de Uso}

\begin{table}[H]
\centering
\begin{tabularx}{\linewidth}{lX}
\toprule
Objetivo & \\
Precondiciones & \\
Post-condiciones & \\
Condición final exitoso & \\
Condición final fallido & \\
Actor primario & \\
Actores secundarios & \\
Secuencia normal & \\
Secuencias alternativas & \\
\bottomrule
\end{tabularx}
\caption{Plantilla de Caso de Uso}
\end{table}

% ──────────────────────────────────────────────────────────
\chapter{Enunciado del Proyecto}
\label{anx:enunciado}

\begin{quote}
Como ya he comentado varias veces en clase, la práctica final consistirá en la elaboración
de una página web que explique las técnicas de usabilidad:

\begin{itemize}
  \item \textbf{Usuario específico}: Desarrolladores de software que conocen la existencia de
    técnicas de usabilidad pero no tienen experiencia aplicándolas.
  \item \textbf{Objetivo específico}: Encontrar técnicas que les permitan tomar decisiones
    centradas en el usuario o evaluar sus desarrollos / prototipos y contar con la información
    suficiente para aprender a ejecutarlas.
  \item \textbf{Contexto de uso}: Deberán ser capaces de consultar esta información de manera
    eficiente desde cualquier tipo de dispositivo capaz de mostrar contenido web. Los usuarios
    pueden usar la web para documentarse, seleccionar y preparar las pruebas desde su puesto de
    oficina, pero también les puede hacer falta consultar la información desde el móvil o una tablet
    mientras están realizando alguna prueba.
\end{itemize}

Para la información de cada técnica que se muestre en la web podréis usar la ficha que hemos
realizado en la última PEC.

\textbf{IMPORTANTE}: El trabajo NO ES SOLO LA WEB, sino que documentéis bien cómo habéis
ido tomando las decisiones de diseño, mediante técnicas de HCI/Usabilidad para llegar a esa web.
ESTO ES ESENCIAL.
\end{quote}

% ──────────────────────────────────────────────────────────
\end{document}
